%1. 基本設定
\documentclass[unicode, aspectratio=169, 12pt]{beamer}
%\usecolortheme[RGB={0, 139, 0}]{structure} %多分RGBで色を変えられる



% 2. 日本語設定（元々はこれだった，明朝体）
% fontspecは指定せず、標準の原ノ味フォント（TeX Live標準）に任せるのが最も安全です
%\usepackage{luatexja} 

% 2. 日本語設定（ゴシック体にフォント変更してみた）
% haranoajiを指定し、スライド全体の日本語デフォルトをゴシック体に変更する
\usepackage[haranoaji]{luatexja-preset}
\renewcommand{\kanjifamilydefault}{\gtdefault}
\renewcommand{\familydefault}{\sfdefault}




% 3. 数式・物理系のパッケージ
\usepackage{amsmath, amssymb, mathtools, bm}
\usepackage{physics} 
% physicsとの競合（\qty）を避ける設定で読み込みます
\usepackage[overwrite-functions=false]{siunitx} 
\usepackage{tensor} % テンソルをいい感じに描くためのパッケージ，それぞれのpcで使えるか要確認！！！！！！！！



% 4. 図・グラフ・枠組み（TikZ & tcolorbox）
\usepackage{graphicx}
\usepackage{tikz}
\usetikzlibrary{math, arrows.meta, angles, quotes, decorations.markings}
\usepackage[most]{tcolorbox} % 枠付きボックス（ascmacの代わりになります）



% 5. レイアウト設定
\renewcommand{\baselinestretch}{1.2} % スライドでは1.5は広すぎるため1.2程度が一般的です
\usepackage{url}

\newcommand{\diff}[2]{\frac{d #1}{d #2}} %常微分を書くパッケージ
\newcommand{\pdiff}[2]{\frac{\mathrm{\partial} #1}{\mathrm{\partial} #2}} %常微分を書くパッケージ

% 6. Beamerのデザインテーマ（Boadillaベース）
\usetheme{Boadilla}
\usecolortheme{whale}
\usefonttheme{professionalfonts}
\setbeamertemplate{navigation symbols}{} % 右下のアイコンを消す
\setbeamerfont{frametitle}{size=\Large, series=\bfseries}
\setbeamersize{text margin left=1cm, text margin right=1cm} %スライドの両端1cmを白紙にして寄りすぎないようにする
\setbeamertemplate{blocks}[rounded][shadow=false] %blockの下側の影を消す
\setbeamertemplate{items}[square] %箇条書きのスタイルを指定
\setbeamertemplate{section in toc}[square] %目次の箇条書きのスタイルを指定
\setbeamertemplate{subsection in toc}[square]

% 7. タイトル情報
\title{最小作用の原理を用いたアインシュタイン方程式の導出}
\author{笹伶夷\and 下橋以紗也\and 田處裕哉}
\institute{大阪大学 理学部 物理学科}
\date{\today}


%===============================================================
% 表記のルールの例
% ベクトル空間は $\bm{V}$
% ベクトルは太字にせずにそのまま x \in \bm{V} のように描く
% テンソルはtensorパッケージを使って  \tensor{t}{^\mu_nu}のように書く (このパッケージじゃないと添字の位置がうまく表示されない)
% 常微分は \diff{f}{t}のように書く, n階微分は \diff{^n f}{t^n}のように書く
% 偏微分は \pdiff{f}{t}のように書く, n階微分は \pdiff{^n f}{t^n}のように書く


%===============================================================
\begin{document}

%タイトルスライド
\begin{frame}
    \titlepage
\end{frame}


%目次スライド
\begin{frame}{目次}
    \tableofcontents
\end{frame}



%===============================================================
\section{特殊相対論の復習と拡張の必要性}


%\begin{frame}{特殊相対論の復習}
%    \begin{block}{慣性系}
%        慣性の法則が成立する観測系のこと.
%
%        特殊相対論では慣性系から見た物理現象を記述する.
%
%        ある慣性系から見た時空の点は$x^\mu = [x^0,x^1,x^2,x^3] =[ct,x,y,z]$で表す.
%    \end{block}
%
%    \begin{block}{ローレンツ変換}
%        ある慣性系$K$に対して等速で動く慣性系$K'$からみた時空の点は
%        \begin{equation*}
%            x'^\mu = \tensor{\Lambda}{^\mu_\nu}x^\nu
%        \end{equation*}
%        で計算できる.
%        この変換では世界間隔と呼ばれる量が不変である.
%        \begin{equation*}
%            ds^2 = \tensor{\eta}{_\mu_\nu}dx^\mu dx^\nu
%        \end{equation*}
%    \end{block}
%\end{frame}

\begin{frame}{特殊相対論の復習}
    \textbf{■ 慣性系と時空の表現}
    \begin{itemize}
        \item 慣性の法則が成立する系で物理現象を記述する
        \item 時空の点は4元位置ベクトル $x^\mu = [ct, x, y, z]^T$ で指定する
    \end{itemize}

    \vspace{0.8em}
    \textbf{■ 慣性系間の変換（ローレンツ変換）}
    \begin{itemize}
        \item 異なる慣性系への座標変換は線形変換となる:
        \begin{equation*}
            x'^\mu = \tensor{\Lambda}{^\mu_\nu} x^\nu
        \end{equation*}
        \item この変換のもとで世界間隔 $ds^2$ は不変に保たれる:
        \begin{equation*}
            ds^2 = \eta_{\mu\nu} dx^\mu dx^\nu = \mathrm{const.} \quad (\eta_{\mu\nu} = \mathrm{diag}(1,-1,-1,-1))
        \end{equation*}
    \end{itemize}
\end{frame}



\begin{frame}{特殊相対論の拡張1:座標変換}
    \begin{block}{特殊相対論で用いる座標}
        直交する4方向の直線を使って時空の点を表現する.

        座標変換は慣性系同士のものに限られている.
    \end{block}

    $\Rightarrow$もっと自由に曲線座標を含めた座標変換をしたい!

    例:
    \begin{align*}
        \left\{
        \begin{aligned}
        x' &= x - \frac{1}{2}at^2\\
        t' &= t
        \end{aligned}
        \right.
        \qquad\qquad
        \left\{
        \begin{aligned}
            r &= \sqrt{x^2 + y^2 + z^2} \\
            \theta &= \arccos\left(\frac{z}{r}\right) \\
            \phi &= \arctan\left(\frac{y}{x}\right)\\
            t' &= t
        \end{aligned}
        \right.
    \end{align*}
\end{frame}

\begin{frame}{特殊相対論の拡張2: 曲がっている空間への対応}
    \textbf{■ 一様な重力場と座標系}
    \begin{itemize}
        \item 一様な重力 $\bm{g}$ が働く系における運動方程式:
        \begin{equation*}
            m \frac{d^2\bm{x}}{dt^2} = \bm{F} + \textcolor{red}{m\bm{g}}
        \end{equation*}
        \item 加速系（自由落下系）への座標変換 $\bm{x}' = \bm{x} - \frac{1}{2}\bm{g}t^2$ を行うと:
        \begin{equation*}
            m \frac{d^2\bm{x}'}{dt^2} = \bm{F} \quad \text{（重力項が消去される）}
        \end{equation*}
    \end{itemize}

    \vspace{0.8em}
    \begin{alertblock}{疑問}
        \centering
        重力は適切な座標変換を選べば\textbf{常に消去できる}のか$\cdots$?
    \end{alertblock}
\end{frame}

\begin{frame}{特殊相対論の拡張2:曲がっている空間への対応}
    \begin{block}{特殊相対論で考える空間}
        直交する4本の直線をどこまでも交わらないように引くことができる「平坦な」空間を考えていた.
    \end{block}

    $\Rightarrow$時空が局所的には平坦でも大局的には「曲がっている」可能性も考慮した理論を作りたい!

    例:
\begin{center}
    \begin{tikzpicture}[scale=0.9, >=latex]
        % --- 左側：曲面（球面） ---
        \begin{scope}[shift={(1.4, -0.5)}]
            % 球体の立体的なグラデーション
            \shade[ball color=blue!20, opacity=0.7] (0,0) circle (1.3);
            
            % 外周の輪郭
            \draw[thick, blue!80!black] (0,0) circle (1.3);
            
            % 緯線（横方向の丸み）
            \draw[blue!60, thin] (-1.3,0) arc (180:360:1.3 and 0.35);
            \draw[blue!30, dashed] (1.3,0) arc (0:180:1.3 and 0.35);
            \draw[blue!60, thin] (-1.13, 0.65) arc (180:360:1.13 and 0.3);
            \draw[blue!60, thin] (-1.13, -0.65) arc (180:360:1.13 and 0.3);
            
            % 経線（縦方向の丸み）
            \draw[blue!60, thin] (0, 1.3) arc (90:270:0.55 and 1.3);
            \draw[blue!60, thin] (0, 1.3) arc (90:-90:0.55 and 1.3);
            
            % 局所的に注目する領域（赤枠）
            \draw[thick, red] (0.15, -0.1) rectangle (0.55, 0.3);
        \end{scope}
        \node[below] at (1.4, -1.85) {\footnotesize 全体（曲がっている）};

        % --- 拡大の点線と矢印 ---
        \draw[dashed, red!60] (1.95, -0.2) -- (4.25, 0.3);
        \draw[dashed, red!60] (1.95, -0.6) -- (4.25, -1.3);
        \draw[->, thick, red] (2.2, -0.5) -- (3.6, -0.5) node[midway, above=-1pt] {\scriptsize 拡大};

        % --- 右側：拡大領域（平らな格子） ---
        \begin{scope}[shift={(5.3, -0.5)}]
            % 円形にクリップして正方形グリッドを描画
            \clip (0,0) circle (1.1);
            \fill[white] (-1.2,-1.2) rectangle (1.2,1.2);
            \draw[step=0.35, blue!70!black, thin] (-1.2,-1.2) grid (1.2,1.2);
        \end{scope}
        \draw[thick, red] (5.3, -0.5) circle (1.1);
        \node[below] at (5.3, -1.85) {\footnotesize 局所的（平ら）};
    \end{tikzpicture}
    \end{center}    
\end{frame}


\section{一般相対論の舞台}

% ここで数学的にローレンツ多様体を導入するのか,時空に対していくつかの条件を要請して作るのかどっちがいいかな

\begin{frame}{一般相対論の舞台}
    \begin{block}{曲がっている時空の表現方法}
        曲がっている時空の点の位置は4個の実数値で指定できる.

        また，点の位置を指定するのに使う実数値の組を\textcolor{red}{局所座標}と呼ぶ.
    \end{block}
    例:

    球面$x^2+y^2+z^2=1$の$z>0$半球上のの点は
    
    緯度と経度を用いて$(\theta,\phi)$で指定することもできるし，
    
    $xy$平面に射影した座標$(x,y)$で指定することもできる.

\end{frame}

%\begin{frame}{一般相対論の舞台}
%    \begin{block}{時空上の物理量}
%        物理量は座標変換$(x^\mu\to y^\mu)$をしたときの変換則で分類される.
%        \begin{itemize}
%            \item スカラー: 座標変換で不変な量
%            \item 反変ベクトル: $A'^\mu = \pdiff{y^\mu}{x^\nu} A^\nu$
%            \item 共変ベクトル: $A'_\mu = \pdiff{x^\nu}{y^\mu} A_\nu$
%            \item テンソル: 複数の添字をもち, 添字ごとに反変・共変の変換を受ける量
%            \begin{itemize}
%                \item 2階反変テンソル: $T'^{\mu\nu} = \pdiff{y^\mu}{x^\rho} \pdiff{y^\nu}{x^\sigma} T^{\rho\sigma}$
%                \item 2階共変テンソル: $T'_{\mu\nu} = \pdiff{x^\rho}{y^\mu} \pdiff{x^\sigma}{y^\nu} T_{\rho\sigma}$
%                \item 2階混合テンソル: $\tensor{{T'}}{^\mu_\nu} = \pdiff{y^\mu}{x^\rho} \pdiff{x^\sigma}{y^\nu} \tensor{T}{^\rho_\sigma}$
%            \end{itemize}
%            \item 一般の$(k, l)$階テンソル $T^{\mu_1 \dots \mu_k}_{\nu_1 \dots \nu_l}$ も各添字について同様に変換する
%        \end{itemize}
%    \end{block}
%\end{frame}


\begin{frame}{一般相対論の舞台}
    \begin{block}{時空上の物理量}
        物理量は座標変換 $(x^\mu \to y^\mu)$ に対する変換則で分類される.
        
        \vspace{0.4em}
        \centering
        \renewcommand{\arraystretch}{1.5}
        \small
        \begin{tabular}{lll}
            \hline
            分類 & 記法 & 変換則 \\
            \hline
            スカラー & $\phi$ & $\phi' = \phi$ \\
            反変ベクトル & $A^\mu$ & $A'^\mu = \pdiff{y^\mu}{x^\nu} A^\nu$ \\
            共変ベクトル & $A_\mu$ & $A'_\mu = \pdiff{x^\nu}{y^\mu} A_\nu$ \\
            2階テンソル & $T^{\mu\nu}$ & $T'^{\mu\nu} = \pdiff{y^\mu}{x^\rho} \pdiff{y^\nu}{x^\sigma} T^{\rho\sigma}$ \\
            \hline
        \end{tabular}
        \vspace{0.3em}
    \end{block}
    \vspace{0.3em}
    $\Rightarrow$ 一般のテンソルは、\textbf{上付き添字1つにつき $\pdiff{y}{x}$}、\textbf{下付き添字1つにつき $\pdiff{x}{y}$} が掛かる.
\end{frame}


\begin{frame}{一般相対論の舞台}
    \begin{block}{時空の2点間の距離の測り方}
        局所座標として$x^\mu$を採用する.

        点P$(x^\mu)$と点Q$(x^\mu + dx^\mu)$の距離の二乗$ds^2$は一般に$dx^\mu$の二次形式
        \begin{equation*}
            ds^2 = g_{\mu_\nu} dx^\mu dx^\nu
        \end{equation*}
        で計算できる.この係数行列$g_{\mu_\nu}$を\textcolor{red}{計量テンソル}と呼ぶ.

        上付きの計量テンソル$g^{\mu_\nu}$を$g_{\mu\nu}$の逆行列として定義する:\>$g_{\mu\nu}g^{\nu\lambda}=\delta^\lambda_\mu$

        これらの計量テンソルを用いてベクトルやテンソルの添字を上下する:
        \begin{align*}
            A^\mu = g^{\mu\nu}A_\nu \quad,\quad A_\mu = g_{\mu\nu}A^\nu
        \end{align*}
    \end{block}
\end{frame}

\begin{frame}{一般相対論の舞台}
    \begin{block}{物理的な要請①}
        宇宙のどこにいても，自由落下する観測者の近くでは慣性系になっている.
    \begin{center}
        $\Downarrow$
    \end{center}
        時空の各点で局所的に計量が$g_{\mu\nu} = \eta_{\mu\nu}$となるような座標を常に選択することができる.
    \end{block}
\end{frame}



\end{document}